% ID: SEL-D0035-OB-E01
% ESTADO: validada. S036, cotejada con examen y criterios.
\textbf{a)} El recinto está por encima de $y=4-2x$ y $y=-1$,
y por debajo de $y=3-x$.
\begin{center}
\setlength{\unitlength}{9mm}
\begin{picture}(8,5)(-1,-2)
\qbezier(0,0)(2.5,0)(5,0)\put(5.1,0){$x$}
\qbezier(0,-1.5)(0,0.5)(0,2.5)\put(0,2.6){$y$}
\qbezier(1,2)(1.75,0.5)(2.5,-1)
\qbezier(2.5,-1)(3.25,-1)(4,-1)
\qbezier(4,-1)(2.5,0.5)(1,2)
\put(1,2.15){\makebox(0,0)[b]{$(1,2)$}}
\put(2.5,-1.2){\makebox(0,0)[t]{$(5/2,-1)$}}
\put(4,-1.2){\makebox(0,0)[t]{$(4,-1)$}}
\end{picture}
\end{center}
\textbf{b)} En $(2,1)$, $x+y=3$, $2x+y=5$ e $y=1$.
Se cumplen las tres restricciones, así que pertenece al recinto.

\textbf{c)} Al igualar las fronteras dos a dos obtenemos
$(1,2)$, $(5/2,-1)$ y $(4,-1)$. En esos vértices $F$ vale
$13$, $17/2$ y $16$. Por linealidad, el mínimo es $\boxed{17/2}$
sólo en $(5/2,-1)$ y el máximo es $\boxed{16}$ sólo en $(4,-1)$.

\textbf{d)} En el lado inferior, $F(x,-1)=5x-4$.
Vale 9 en $x=13/5$, que está entre $5/2$ y 4.
Por tanto $\boxed{(13/5,-1)}$ es un punto factible con $F=9$.
